\documentclass[12pt, a4paper]{article}

\usepackage[utf8]{inputenc}
\usepackage[T1]{fontenc}
\usepackage[italian]{babel}
\usepackage{amsmath}
\usepackage{amssymb}
\usepackage{geometry}
\usepackage{enumitem}
\usepackage{pgfplots}
\pgfplotsset{compat=1.18}
\usepackage{tikz}
\usepackage{xcolor}
\usepackage{booktabs}
\usepackage{array}
\usepackage{parskip}

\geometry{margin=2.5cm}

\definecolor{gaussblue}{RGB}{30, 90, 160}
\definecolor{lightblue}{RGB}{220, 235, 255}
\definecolor{darkgray}{RGB}{60, 60, 60}

\title{%
  \textcolor{gaussblue}{\textbf{Distribuzione gaussiana e deviazione standard}}
}
\author{Prof. Fedeli Massimo}
\date{IIS Fermi Sacconi Cpia di Ascoli Piceno}

\begin{document}

\maketitle
\thispagestyle{empty}

% ─────────────────────────────────────────────
\section{Distribuzione gaussiana}

La \textbf{distribuzione gaussiana}, detta anche \textbf{distribuzione normale}, è uno dei modelli statistici più utilizzati per descrivere fenomeni reali. 

Il grafico di questa distribuzione ha una caratteristica \emph{forma a campana}. I valori più \textbf{frequenti} si trovano vicino al centro della distribuzione, mentre valori molto piccoli o molto grandi diventano progressivamente meno probabili.

\begin{center}
\begin{tikzpicture}
\begin{axis}[
  width=10cm, height=6cm,
  xlabel={$x$},
  ylabel={$f(x)$},
  axis lines=left,
  xtick={-3,-2,-1,0,1,2,3},
  xticklabels={$\mu{-}3\sigma$, $\mu{-}2\sigma$, $\mu{-}\sigma$, $\mu$,
               $\mu{+}\sigma$, $\mu{+}2\sigma$, $\mu{+}3\sigma$},
  ytick=\empty,
  domain=-3.5:3.5,
  samples=200,
  smooth,
  thick,
  color=gaussblue,
  fill=lightblue,
  fill opacity=0.5,
]
\addplot [gaussblue, thick, fill=lightblue, fill opacity=0.5]
  {exp(-x^2/2) / sqrt(2*pi)} \closedcycle;
\end{axis}
\end{tikzpicture}
\end{center}

\underline{Molti fenomeni naturali seguono approssimativamente questo modello.} Alcuni esempi sono:

\begin{itemize}
  \item altezza delle persone in una popolazione;
  \item errori nelle misurazioni scientifiche;
  \item rumore nei segnali elettronici;
  \item risultati di test statistici.
\end{itemize}

La distribuzione gaussiana è descritta da due parametri principali:

\begin{itemize}
  \item \textbf{media} ($\mu$): rappresenta il valore centrale della distribuzione;
  \item \textbf{deviazione standard} ($\sigma$): misura quanto i valori sono dispersi rispetto alla media.
\end{itemize}

La distribuzione è \emph{simmetrica} rispetto alla media: valori sopra e sotto la media hanno la stessa probabilità.

% ─────────────────────────────────────────────
\section{Deviazione standard}

La \textbf{deviazione standard} è una misura statistica che indica quanto i dati sono dispersi rispetto alla media.

\begin{itemize}
  \item Se i dati sono molto vicini alla media, la deviazione standard è \emph{piccola}.
  \item Se i dati sono molto lontani dalla media, la deviazione standard è \emph{grande}.
\end{itemize}

\subsection{Formula della deviazione standard}

Per una popolazione di $N$ valori $x_1, x_2, \ldots, x_N$ con media $\mu$, la deviazione standard è definita come:

\[
  \sigma = \sqrt{\dfrac{1}{N} \sum_{i=1}^{N} (x_i - \mu)^2}
\]

Il calcolo avviene in tre passaggi logici:

\begin{enumerate}
  \item Si calcola la \textbf{media} dei dati.
  \item Si calcola la \textbf{distanza} di ogni valore dalla media.
  \item Si fa la media dei \textbf{quadrati} di queste distanze e si calcola la radice quadrata.
\end{enumerate}

Il quadrato serve a evitare che le differenze positive e negative si annullino tra loro.

\subsection{Esempio intuitivo}

Consideriamo due gruppi di dati.

\begin{center}
\begin{tabular}{lcc}
\toprule
\textbf{Gruppo} & \textbf{Valori} & \textbf{Media} \\
\midrule
A & $10,\ 11,\ 9,\ 10,\ 10$ & $10$ \\
B & $2,\ 18,\ 5,\ 15,\ 10$  & $10$ \\
\bottomrule
\end{tabular}
\end{center}

\medskip
\noindent\textbf{Calcolo per il Gruppo A} \quad ($\mu = 10$)

\begin{center}
\begin{tabular}{ccc}
\toprule
$x_i$ & $x_i - \mu$ & $(x_i - \mu)^2$ \\
\midrule
$10$ & $0$  & $0$ \\
$11$ & $+1$ & $1$ \\
$9$  & $-1$ & $1$ \\
$10$ & $0$  & $0$ \\
$10$ & $0$  & $0$ \\
\midrule
     & \textbf{Somma} & $2$ \\
\bottomrule
\end{tabular}
\end{center}

\[
  \sigma_A = \sqrt{\dfrac{2}{5}} = \sqrt{0{,}4} \approx \mathbf{0{,}63}
\]

\medskip
\noindent\textbf{Calcolo per il Gruppo B} \quad ($\mu = 10$)

\begin{center}
\begin{tabular}{ccc}
\toprule
$x_i$ & $x_i - \mu$ & $(x_i - \mu)^2$ \\
\midrule
$2$  & $-8$ & $64$ \\
$18$ & $+8$ & $64$ \\
$5$  & $-5$ & $25$ \\
$15$ & $+5$ & $25$ \\
$10$ & $0$  & $0$  \\
\midrule
     & \textbf{Somma} & $178$ \\
\bottomrule
\end{tabular}
\end{center}

\[
  \sigma_B = \sqrt{\dfrac{178}{5}} = \sqrt{35{,}6} \approx \mathbf{5{,}97}
\]

\medskip
\noindent\textbf{Confronto:}
\begin{itemize}
  \item \textbf{Gruppo A}: $\sigma_A \approx 0{,}63$ — i valori sono molto vicini alla media, la deviazione standard è \emph{piccola}.
  \item \textbf{Gruppo B}: $\sigma_B \approx 5{,}97$ — i valori sono molto dispersi, la deviazione standard è \emph{grande}.
\end{itemize}

Pur avendo la stessa media ($\mu = 10$), i due gruppi hanno distribuzioni molto diverse: la deviazione standard cattura esattamente questa differenza.

% ─────────────────────────────────────────────
\section{Regola 68–95–99{,}7}

Una proprietà importante della distribuzione normale è la cosiddetta \textbf{regola empirica}:

\begin{itemize}
  \item circa \textbf{68\%} dei valori si trova nell'intervallo $\mu \pm \sigma$;
  \item circa \textbf{95\%} dei valori si trova nell'intervallo $\mu \pm 2\sigma$;
  \item circa \textbf{99{,}7\%} dei valori si trova nell'intervallo $\mu \pm 3\sigma$.
\end{itemize}

\begin{center}
\begin{tikzpicture}
\begin{axis}[
  width=11cm, height=6cm,
  axis lines=left,
  xtick={-3,-2,-1,0,1,2,3},
  xticklabels={$\mu{-}3\sigma$, $\mu{-}2\sigma$, $\mu{-}\sigma$, $\mu$,
               $\mu{+}\sigma$, $\mu{+}2\sigma$, $\mu{+}3\sigma$},
  ytick=\empty,
  domain=-3.5:3.5,
  samples=200,
  smooth,
]
% 99.7% band
\addplot [fill=blue!15, draw=none, domain=-3:3] {exp(-x^2/2)/sqrt(2*pi)} \closedcycle;
% 95% band
\addplot [fill=blue!30, draw=none, domain=-2:2] {exp(-x^2/2)/sqrt(2*pi)} \closedcycle;
% 68% band
\addplot [fill=blue!55, draw=none, domain=-1:1] {exp(-x^2/2)/sqrt(2*pi)} \closedcycle;
% curve
\addplot [gaussblue, thick, domain=-3.5:3.5] {exp(-x^2/2)/sqrt(2*pi)};
% labels
\node at (axis cs:0,0.18) [font=\small\bfseries, white] {68\%};
\node at (axis cs:1.5,0.07) [font=\small, gaussblue] {95\%};
\node at (axis cs:2.6,0.025) [font=\small, darkgray] {99{,}7\%};
\end{axis}
\end{tikzpicture}
\end{center}

Questo significa che la maggior parte dei dati è concentrata vicino alla media.

% ─────────────────────────────────────────────
\section{Conclusione}

La distribuzione gaussiana è un modello fondamentale in statistica. La \textbf{media} indica il centro della distribuzione, mentre la \textbf{deviazione standard} descrive quanto i dati si disperdono attorno a questo valore. Comprendere questi due concetti permette di interpretare correttamente molti fenomeni osservati nel mondo reale.

\end{document}
